\documentclass[11pt]{article}
\usepackage{hyperref} % for referencing
\usepackage{graphicx} % Required for inserting images
\usepackage{fancyhdr} % for the page numbering, footers, headers, etc.
\usepackage{geometry} % for the page structure
\usepackage{setspace} % for line spacing
\usepackage{bm} % for easy vectors

% page structure
\geometry{a4paper, total={170mm, 257mm}, left=30mm, right=20mm, top=25mm, bottom=25mm}

% commands the the citations I am using
\renewcommand{\cite}[1]{[\ref{#1}]} 
%\renewcommand{\labelenumi}{[\theenumi]}

% commands for the page number
\pagestyle{fancy}
\fancyhf{}
\fancyfoot[R]{\thepage}
\renewcommand{\headrulewidth}{0pt}

% commands for line spacing
\onehalfspacing

% no paragraph indent
\setlength{\parindent}{0pt}

\begin{document}

\thispagestyle{empty} % no specific page style here

Title page \\
Test

\newpage

\setcounter{page}{1} % sets the page counter to 1

\tableofcontents

\newpage

\section*{Abstract}

\newpage

\section{Introduction and Motivation}

- Lymph Nodes and the Immune System \\
- Structure of the Lymph Nodes \\

- Agent Based Models in cell and tissue mechanics \\

(https://pmc.ncbi.nlm.nih.gov/articles/PMC7063836/)
(https://pmc.ncbi.nlm.nih.gov/articles/PMC6557636/)

\newpage

\section{Methods}

We now want to introduce the mathematical and physical methods required to dynamically model the structure of the fibroblastic reticular network in the lymph nodes. After having established these methods we will further introduce the ya||a (yet another parallel agent based model for morphogenesis) framework which will be the main simulation tool we will use. \par

\subsection{Physically Modelling Cells}

In a very simplified world cells can be viewed as particles in an environment following certain rules that dictate their dynamics. Such a set-up is very common in physics, most first physics courses start with describing the dynamics of such particles. Therefore, if we decide such a view on cells we will be able to use the vast pool of methods that classical mechanics has to offer. There are two main ways in classical mechanics to describe systems made of particles: point-like theories and field theories.

\subsubsection*{Point-Like Theories}

In point-like dynamics every particle's properties, e.g. position, velocity, etc., are describes by their own function. We can therefore characterize the entire system by a set of functions: 

\begin{equation}
    \label{point-like-set}
    \mathcal{P}_t = \{\bm{f}_1(t), ..., \bm{f}_N(t), \bm{g}_1(t), ... \bm{g}_N(t),...\}
\end{equation}

To describe the dynamics of such a system are then described by a set of ordinary differential equations in time.

\begin{equation}
    \label{point-like-ode}
    \frac{d}{dt} \bm{f}_i (t) = \bm{F}_i (\mathcal{P}_t, t)
\end{equation}

Here $\bm{F}_i(\mathcal{P}_t, t)$ describes the dynamics of the the ith particle in time and due to the other particles in the system. \par

\vspace{4mm}

In physical systems the quantities of interest are usually the position and velocity of the particles, as they allow us to derive other quantities of interest, e.g. the kinetic energy, potential energy, etc.. Furthermore, the ordinary differential equations that describe the dynamics of the system are often determined by the specific environment and particles in question through the use of physical laws and fundamental principles such as Newton's Second Law of Motion ($F = m \bm{a}$), the Euler-Lagrange Equation ($\frac{d}{dt} \frac{\partial L}{\partial \dot{q}} = \frac{\partial L}{\partial q}$) or the Hamiltonian Equations of Motion ($\frac{\partial H}{\partial p} = \dot{q}, \ \frac{\partial H}{\partial \dot{q}} = - \dot{p}$).

\subsubsection*{Field Theories}

Field theories, rather than describing every particle seperately, describe the dynamics of multi-variable functions of space and time. Instead of describing the position of every particle in space we would now describe the density of particles throughout all of space, so one function in space and time instead of N function of time. This works equivalently for tracking vector valued functions such as velocity. 

\newpage

The entire system can therefore be expressed in terms of a smaller set of variables: 


\begin{equation}
    \label{field-set}
    \mathcal{F}_t = \{\rho(\bm{x}, t), \bm{v}(\bm{x}, t), ... \}
\end{equation}

The dynamics of these fields are then described by a set of partial differential equations:


\begin{equation}
    \label{field-pde}
    \frac{d \bm{f}}{dt} = \hat{D} (\mathcal{F}_t, t)
\end{equation}

Here $\hat{D}$ is a differential opertor acting on the fields of the system and time. \par

\vspace{4mm}

Identically to before the dynamics of the field in physics is determined by the spefics of the system itself, physical laws and fundamental principles like the Euler-Lagrange Equation for fields ($\partial_\mu \left( \frac{\partial \mathcal{L}}{\partial (\partial_\mu q)} \right)= \frac{\partial \mathcal{L}}{\partial q}$).

\vspace{4mm}

Both of these approaches to modelling systems of particle are helpful and find use all throughout physics, particularly field theory is the standard in modern physics with the success of the Standard Model through Quantum Field Theory and General Relativity. \par

\subsubsection*{Modelling Cells on the Computer}

As we are interested in simulating cells on a computer it is a lot easier for us to work with the point-like description. We will therefore now layout how exactly we will model cells and their dynamics on the computer. \par

\vspace{4mm}



\newpage

\section{Data Analysis}

\newpage

\section{Outlook}

\renewcommand{\labelenumi}{[\theenumi]}

\section*{References}
\begin{enumerate}
    \item nothing
\end{enumerate}

\end{document}

